\chapter{Introducción}

Las redes satelitales, particularmente aquellas compuestas por constelaciones de órbita baja (LEO), representan un desafío fundamental para los protocolos de comunicaciones tradicionales debido a su topología altamente dinámica y la intermitencia inherente de sus enlaces. En estos entornos, caracterizados por latencias de propagación variables y ventanas de visibilidad limitadas, resulta indispensable el uso de arquitecturas de redes tolerantes a demoras o interrupciones (DTN).

En la actualidad, el protocolo de enrutamiento más extendido para este tipo de escenarios es \textit{Contact Graph Routing} (CGR) \cite{burleigh-dtnrg-cgr-01}. CGR basa sus decisiones en el cálculo de rutas óptimas a partir de un conocimiento previo y preciso de los contactos futuros entre los nodos de la red. Si bien este enfoque es efectivo en misiones con movilidad determinista, presenta limitaciones críticas en términos de escalabilidad y adaptabilidad. La necesidad de contar con una visión global y actualizada del plan de contactos (\textit{contact plan}) dificulta su respuesta ante cambios imprevistos en la topología o fallas inesperadas en los nodos, volviéndose limitante en escenarios de gran complejidad.

Frente a estas limitaciones, este trabajo explora la viabilidad del protocolo \textit{Adjacent Network Topology} (ANTop) \cite{marcu2007antop} como una solución eficiente y adaptable para el enrutamiento satelital. A diferencia de los esquemas que dependen de información de estado global, ANTop codifica la topología de la red directamente en las direcciones de los nodos. Esto habilita un esquema de ruteo reactivo donde cada satélite puede tomar decisiones de reenvío de forma puramente local, utilizando únicamente información de sus vecinos inmediatos.

El objetivo principal de este proyecto es la implementación de ANTop en un entorno de simulación de redes satelitales, integrando el sistema de indexación geoespacial H3 para representar la superficie terrestre. A través de este desarrollo, se busca realizar una comparación de desempeño frente a CGR, analizando métricas de latencia, tasa de entrega y costo computacional. De esta manera, se pretende demostrar que un enfoque descentralizado y basado en la posición relativa de los nodos puede ofrecer una mayor robustez y escalabilidad en las constelaciones satelitales modernas.
